\section{Summation and passage to measurable inputs}\label{sec:summation}

The local increase from Proposition~\ref{prop:increase} has the sign
needed for telescoping toward finer scales. This use of an energy
difference parallels the Bellman-function summations in
\cite[Section~2]{Kovac2011} and \cite[Lemma~2]{Kovac2012}.
We bound the total energy at each level before carrying out that sum.

\begin{lemma}[Total energy]\label{lem:total}
For the step functions of Section~\ref{sec:dyadic} and every \(k\ge m\),
set
\[
 J_k=\sum_{\mathbf I\in\cA_k}2^{2k}e(\mathbf I).
\]
Then
\begin{equation}\label{eq:total}
 0\le J_k\le2\sqrt2\sum_{v=0}^2\norm{F_v}_3^3.
\end{equation}
For \(k\ge m+1\),
\begin{equation}\label{eq:telescope}
 \sum_{\mathbf I\in\cA_k}2^{2k}\delta(\mathbf I)
       =J_{k-1}-J_k.
\end{equation}
\end{lemma}

\begin{proof}
If \(a,b\) are the Hilbert--Schmidt norms of the two edges incident
to \(v\), then
\[
 e_v(\mathbf I)
 =\tr(T_v^{3/2})
 \le(\tr T_v)^{3/2}
 =(a^2+b^2)^{3/2}
 \le\sqrt2(a^3+b^3).
\]
For an edge \(W_v\), its probability normalization gives
\begin{equation}\label{eq:edge-mass}
 \norm{W_v}_{\HS}^3
 =\left(\frac1{\abs{I_v}\abs{I_{v+1}}}
          \int_{I_v\times I_{v+1}}\abs{F_v}^2\right)^{3/2}
 \le \frac1{\abs{I_v}\abs{I_{v+1}}}
          \int_{I_v\times I_{v+1}}\abs{F_v}^3.
\end{equation}
Every edge occurs in two row stars. Also, every ordered pair of
intervals in \(\cD_k\) determines exactly one admissible third
interval, by XOR. Multiplication by \(2^{2k}=\abs{I_v}\abs{I_{v+1}}\)
and summation of \eqref{eq:edge-mass} therefore prove \eqref{eq:total}.
These sums are finite: compact support leaves finitely many nonzero
edge matrices, each belonging to one admissible triple.

Deleting the last bit of each child index shows that every admissible
child has a unique admissible parent. As
\(2^{2k}/4=2^{2(k-1)}\), the child part in the definition of
\(\delta\) sums to \(J_{k-1}\), proving \eqref{eq:telescope}.
\end{proof}

\begin{proof}[Proof of Theorem~\ref{thm:main} for step functions]
The case \(S=\varnothing\) is immediate. Otherwise set
\(k_-=\min S\), \(k_+=\max S\), and refine the common atom level,
if necessary, so that \(m\le k_--1\). By \eqref{eq:trace-integral},
Proposition~\ref{prop:increase}, and Lemma~\ref{lem:total},
\begin{align*}
 \sum_{k\in S}\sum_{\mathbf I\in\cA_k}\abs{L_{\mathbf I}(F_0,F_1,F_2)}
 &\le\frac{10\sqrt2}{3}
       \sum_{k=k_-}^{k_+}(J_{k-1}-J_k)\\
 &\le\frac{10\sqrt2}{3}J_{k_--1}
 \le\frac{40}{3}\sum_{v=0}^2\norm{F_v}_3^3.
\end{align*}
The insertion of missing scales is valid because each energy
difference is nonnegative. If every input has nonzero norm, apply
this estimate to \(F_v/\norm{F_v}_3\). Trilinearity of each local
integral gives \eqref{eq:main}, with constant \(40\).
If any norm is zero, the local integrals vanish.
\end{proof}

To pass to general inputs we need continuity of the absolute sum,
not merely continuity of the signed form. The following fixed-level
estimate provides it.

\begin{lemma}[Continuity at one scale]\label{lem:continuity}
For every \(k\in\Z\) and real \(F_v\in L^3(\R_+^2)\),
\begin{equation}\label{eq:one-scale}
 \sum_{\mathbf I\in\cA_k}\abs{L_{\mathbf I}(F_0,F_1,F_2)}
       \le\prod_{v=0}^2\norm{F_v}_3.
\end{equation}
Consequently the sequence
\((L_{\mathbf I}(F_0,F_1,F_2))_{\mathbf I\in\cA_k}\) depends
continuously on the inputs as a map from \(L^3\times L^3\times L^3\)
to \(\ell^1(\cA_k)\).
\end{lemma}

\begin{proof}
Define the nonnegative edge averages
\[
 a_v(\mathbf I)=
 \frac1{\abs{I_v}\abs{I_{v+1}}}
            \int_{I_v\times I_{v+1}}\abs{F_v}^3.
\]
H\"older's inequality on the product probability space of the three
intervals, with the Haar signs bounded by one, gives
\[
 \abs{L_{\mathbf I}(F_0,F_1,F_2)}
       \le2^{2k}\prod_{v=0}^2 a_v(\mathbf I)^{1/3}.
\]
Apply H\"older again, now to the sum over \(\cA_k\) with weight
\(2^{2k}\). The uniqueness of the third interval yields
\[
 \sum_{\mathbf I\in\cA_k}2^{2k}a_v(\mathbf I)=\norm{F_v}_3^3,
\]
proving \eqref{eq:one-scale}. These bounds also justify all integrals
and absolute sums for \(L^3\) inputs.

For two input triples \(F,G\), expand the difference of local
integrals into the three terms with inputs
\[
 (F_0-G_0,F_1,F_2),\quad
 (G_0,F_1-G_1,F_2),\quad
 (G_0,G_1,F_2-G_2).
\]
Estimate each term by \eqref{eq:one-scale}. If \(G_v\to F_v\) in
\(L^3\), the sum of these three product bounds tends to zero.
This is the asserted \(\ell^1\) continuity.
\end{proof}

\begin{proof}[Completion of Theorem~\ref{thm:main}]
Bounded compactly supported dyadic step functions are dense in
\(L^3(\R_+^2)\). To see this, first approximate an \(L^3\) function by
a continuous compactly supported function, then take its averages on
sufficiently small dyadic squares. Uniform continuity and a common
bounded support give convergence in \(L^3\).

For arbitrary real \(F_v\in L^3\), choose such step functions
\(F_v^{(n)}\to F_v\) in \(L^3\). Assuming \(S\ne\varnothing\),
refine their common atom levels, without changing the functions, to
satisfy \(m_n\le\min S-1\). Apply the step-function estimate to these
inputs. Lemma~\ref{lem:continuity}, summed over the fixed finite set
\(S\), gives convergence of their local integrals in \(\ell^1\).
The inequality \(\abs{\abs a-\abs b}\le\abs{a-b}\) permits passage to
the absolute sum in \eqref{eq:main}, and norm convergence handles its
right-hand side. This proves the theorem and its stated extension
to real \(L^3\) inputs. The case \(S=\varnothing\) remains immediate.

For any coefficients
\(\abs{\epsilon_{\mathbf I}}\le1\), the triangle inequality gives
the corresponding bound for
\[
 \sum_{k\in S}\sum_{\mathbf I\in\cA_k}
 \epsilon_{\mathbf I}L_{\mathbf I}(F_0,F_1,F_2).
 \qedhere
\]
\end{proof}
